% chapters/07_projektverlauf.tex

\chapter{Projektverlauf und Lessons Learned}
\label{chap:projektverlauf}

Die fachlichen Befunde sind in \cref{chap:ergebnisse,chap:diskussion}
dargestellt. Dieses Kapitel wechselt die Perspektive und reflektiert den
Projekt\emph{verlauf}: die größeren Hindernisse und wie mit ihnen umgegangen
wurde (\cref{sec:hindernisse}), die daraus abgeleiteten Lehren für ein
nächstes vergleichbares Projekt (\cref{sec:lessons}), die Wahl der
Arbeitsmittel (\cref{sec:arbeitsmittel}) und die Rolle KI-gestützter
Entwicklung (\cref{sec:ki-entwicklung}). Das Fazit in \cref{chap:fazit}
bleibt davon getrennt fachlich.

\section{Hindernisse im Projektverlauf}
\label{sec:hindernisse}

Organisatorisch prägend war die \textbf{Beschaffungsphase am Projektbeginn}:
Bevor der erste Trainingslauf starten konnte, mussten die Rechenressourcen
überhaupt erst verfügbar gemacht werden -- der Zugang zum KISSKI-Cluster,
die Evaluation von Cloud-GPU-Anbietern (vast.ai u.\,a.) als Überbrückung und
die Abstimmung über den institutseigenen GPU-Server. Das kostete zu Beginn
erhebliche Kalenderzeit, die nicht durch Mehrarbeit kompensierbar war, weil
sie aus Antrags- und Wartefristen bestand. Ein zweiter organisatorischer
Engpass zog sich länger hin: die \textbf{Wartezeit auf den Zugriff auf den
institutseigenen Server}. Die Closed-Loop-Simulation wich deshalb zunächst
auf gemietete vast.ai-Instanzen aus; erst mit dem Serverzugang (ab Juli~2026)
wurden Referenz-Eval und RL-Fine-Tuning praktikabel, da beide lange
interaktive Sitzungen auf RT-Core-Hardware erfordern.
\Cref{tab:hindernisse} fasst die größeren Hindernisse zusammen; die
technischen sind an den genannten Stellen im Detail behandelt.

\begin{table}[htb]
  \centering
  \begin{tabularx}{\linewidth}{>{\raggedright\arraybackslash}p{5.6cm}X}
    \toprule
    \textbf{Hindernis} & \textbf{Auswirkung und Umgang} \\
    \midrule
    Beschaffung der Rechenressourcen zum Projektstart
      & verzögerter Beginn der Trainingsläufe; KISSKI-Antrag, vast.ai als
        Überbrückung \\
    Wartezeit auf den institutseigenen GPU-Server
      & Sim-Eval nur auf Miet-GPUs, RL unmöglich; ab Juli~2026
        Blackwell-Server \\
    A100/H100 ohne RT-Cores (\cref{sec:eval-methodik})
      & Cluster kann die Sim-Eval nicht ausführen; Aufgabenteilung Training
        (KISSKI) / Simulation (RT-Core-Hardware) \\
    Visueller Domain-Gap (\cref{sec:diagnose,sec:span-gate})
      & Closed-Loop-Sim als Benchmark unbrauchbar ($0/20$); diagnostiziert,
        Gegenmaßnahmen: Visual Matching, Co-Training, RL \\
    Isaac-Sim-Absturz auf Blackwell (\cref{sec:neukalibrierung})
      & Sim-Stack blockiert; Migration auf Isaac~Sim~6.0 \\
    Kamera-Orientierungsbug 95--137° (\cref{sec:neukalibrierung})
      & fünf Diagnose-Läufe fehlinterpretiert; unabhängige USD-Messquelle,
        berechnete Neukalibrierung \\
    Widerrufener Vorzeichen-Fix (\cref{sec:vorzeichen-widerruf})
      & alle Sim-Greifzahlen unter Vorbehalt; Rücknahme, Neumessung
        ausstehend \\
    Host-RAM-Engpass Multi-GPU (\cref{sec:multi-gpu})
      & abgebrochene Jobs; weniger Dataloader-Worker, mehr Speicher \\
    Fork erzwingt \texttt{eval\_strategy="no"} (\cref{sec:eval-methodik})
      & keine Validierung im Training; nachgelagerter Checkpoint-Sweep \\
    \bottomrule
  \end{tabularx}
  \caption{Die größeren Hindernisse des Projekts: organisatorische (oben)
    und technische (unten) mit Auswirkung und Umgang.}
  \label{tab:hindernisse}
\end{table}

\section{Lessons Learned}
\label{sec:lessons}

Aus dem Verlauf lassen sich konkrete Lehren ableiten -- formuliert als das,
was in einem nächsten vergleichbaren Projekt \emph{anders} liefe:

\begin{description}
  \item[Validierung von Beginn an erzwingen.] Held-out-Split und
    Checkpoint-Sweep kamen erst mit dem dritten Lauf -- die Läufe~1 und~2
    werteten dadurch blind den letzten Checkpoint aus, der nachweislich um
    25\,\% schlechter ist als der beste (\cref{sec:lauf3}). Künftig gehören
    Split und nachgelagerte Checkpoint-Auswahl zur Definition-of-Done des
    \emph{ersten} Laufs, nicht zur Nachrüstung.
  \item[Das Messinstrument vor dem Experiment validieren.] Die
    SIMPLER-Literatur hätte vorhergesagt, dass Closed-Loop-Sim ohne Visual
    Matching für dieses Modell $\approx 0\,\%$ liefert
    (\cref{sec:einordnung-null}) -- gelesen wurde sie erst nach dem
    $0/20$-Befund. Vor dem Training klären: Womit wird Erfolg gemessen, und
    ist dieses Instrument für das konkrete Modell überhaupt valide?
  \item[Referenz-Benchmarks früh fahren.] Die RoboCasa-Gegenprobe schloss
    die eigene Inferenz-Kette als Fehlerquelle aus (\cref{sec:robocasa}) --
    aber erst im Juli. Früher gefahren hätte sie Wochen an Unsicherheit
    gespart. Verallgemeinert: Jede selbstgebaute Eval-Kette zuerst an einem
    publizierten Sollwert reproduzieren, dann eigene Messungen vertrauen.
  \item[Grundwahrheit vor abgeleiteten Messungen beschaffen.] Szene und
    Geometrie wurden erst spät gegen die Bilddaten verifiziert; drei
    unabhängige Geometriefehler und der fehlerhafte Vorzeichen-Fix blieben
    dadurch monatelang unentdeckt (\cref{sec:geometrie}). Die Messregeln aus
    \cref{sec:messartefakte} (unabhängige Kontrollgrößen, Kontaktflächen
    statt Körper-Frames, Sättigungszustände protokollieren) gehören von
    Projektbeginn an institutionalisiert, nicht als Reaktion auf Artefakte.
  \item[Hardware-Anforderungen als Matrix, nicht als Summe.] Dass die
    stärksten Trainings-GPUs (A100/H100) die Simulation \emph{nicht}
    ausführen können, wurde erst beim Aufsetzen der Sim-Eval entdeckt.
    Anforderungen wie RT-Cores, Treiber-/Simulator-Kompatibilität und
    VRAM-Klassen sind vorab pro Arbeitspaket zu erheben
    (\cref{sec:arbeitsmittel}).
  \item[Dokumentation als Projektgedächtnis führen.] Diagnose-Chronik,
    TL;DR-Konvention und automatische Drift-Checks machten Widerrufe wie den
    Vorzeichen-Fix überhaupt erst sauber rückverfolgbar
    (\doc{docs/ergebnisse/diagnose-chronik.md}). Diese Disziplin früh zu
    etablieren kostet wenig und zahlt sich spätestens bei der ersten
    zurückgenommenen \enquote{Korrektur} aus.
\end{description}

\section{Wahl der Arbeitsmittel}
\label{sec:arbeitsmittel}

Die wichtigste organisatorische Empfehlung an vergleichbare Projekte lautet:
\textbf{Rechenressourcen vor Projektbeginn klären und beantragen} --
Antrags- und Wartefristen bestimmen sonst den kritischen Pfad, wie die
Beschaffungsphase dieses Projekts gezeigt hat (\cref{sec:hindernisse}).

Inhaltlich hat sich dabei die \textbf{größere, professionell betriebene
Infrastruktur} (KISSKI-Klasse) gegenüber einem kleinen Institutsserver als
vorzugswürdig erwiesen: bei der \emph{Kapazität} (A100/H100 mit
80--94~GB VRAM; der Vier-GPU-Lauf absolvierte dieselbe Datenmenge in 8,4
statt 21,4~Stunden), bei der \emph{Verfügbarkeit} (\ac{slurm}-Warteschlange und
professioneller Betrieb statt Konkurrenz um einen Einzelserver samt eigener
Treiber- und Wartungsverantwortung) und bei der \emph{Skalierbarkeit}
(Batch-Größen und parallele Läufe, die lokal unerreichbar sind). Zugleich
zeigt der Verlauf, dass ein Cluster allein nicht genügt: Sonderanforderungen
wie RT-Core-Rendering, lange interaktive Diagnose-Sitzungen und das
RL-Fine-Tuning brauchten den lokalen Blackwell-Server. Die Konsequenz ist
keine Entweder-oder-Entscheidung, sondern eine gezielte Kombination -- das
skalierende Training auf den Cluster, Simulation und Interaktives auf
dedizierte lokale Hardware -- auf Basis einer vorab erstellten
Anforderungsmatrix (\cref{sec:lessons}). Einzuplanen ist außerdem, dass sehr
neue Hardware Software-Risiken mitbringt: Auf der Blackwell-Generation lief
der etablierte Simulator-Stack erst nach Migration auf Isaac~Sim~6.0
(\cref{sec:neukalibrierung}).

\section{KI-gestützte Entwicklung}
\label{sec:ki-entwicklung}

Projektbegleitend wurde mit \textbf{Claude Code} ein agentisches
KI-Entwicklungswerkzeug eingesetzt (inklusive projektspezifischer Skills,
MCP-Server-Anbindungen u.\,a.\ an Weights\,\&\,Biases und GitHub sowie
Subagenten für Recherche- und Prüfaufgaben). Für die in dieser Arbeit
erreichte Breite -- Trainingsinfrastruktur, Simulationsumgebung, die
Diagnose-Werkzeugkette aus \cref{tab:werkzeuge} und eine umfangreiche
technische Dokumentation neben dem eigentlichen Studienbetrieb -- war diese
Arbeitsweise praktisch unerlässlich. Konkret trug sie bei: bei der
Entwicklung und Iteration der Diagnoseskripte, bei der statistischen
Auswertung der Läufe und der Pflege der Ergebnisdokumente, bei größeren
Refactorings (Bedienoberfläche, Portabilität) und bei der
Konsistenzsicherung der Dokumentation (Konventions- und Drift-Prüfungen).

Einzuordnen ist das nüchtern: Das Werkzeug ersetzt weder Fachurteil noch
Verifikation. Die Messartefakt-Lehren aus \cref{sec:messartefakte} gelten
unverändert auch für KI-generierte Diagnosen -- jeder maßgebliche Befund
wurde gegen unabhängige Messungen geprüft, und Fehlschlüsse (etwa der
zunächst plausible Vorzeichen-Fix) entstanden und fielen unabhängig vom
Werkzeug durch dieselbe methodische Kontrolle. Der belastbare Effekt war
eine deutlich höhere Umsetzungs- und Diagnosegeschwindigkeit bei
gleichbleibender Verantwortung der Verfasser für alle Inhalte; die formale
Offenlegung erfolgt in der Erklärung zur KI-Nutzung am Ende der Arbeit.
